% Part: first-order-logic
% Chapter: sequent-calculus
% Section: quantifier-rules

\documentclass[../../../include/open-logic-section]{subfiles}

\begin{document}

\olfileid[hi]{fol}{seq}{qrl}

\olsection{परिमाणक-नियम}

\subsection{$\lforall$ के नियम}

\begin{defish}
\Axiom$ !A(t), \Gamma \fCenter \Delta$
\RightLabel{\LeftR{\lforall}}
\UnaryInf$ \lforall[x][!A(x)],\Gamma \fCenter \Delta$
\DisplayProof
\hfill
\Axiom$ \Gamma \fCenter \Delta, !A(a) $
\RightLabel{\RightR{\lforall}}
\UnaryInf$ \Gamma \fCenter \Delta, \lforall[x][!A(x)]$
\DisplayProof
\end{defish}

\LeftR{\lforall} में $t$ बंद पद है (अर्थात् उसमें कोई चर नहीं है)।
\RightR{\lforall} में $a$~ऐसा !!a{constant} है जो \RightR{\lforall} नियम के
निचले अनुक्रम में कहीं भी न आए। $a$ को \RightR{\forall} अनुमान का
\emph{अभिलाक्षणिक चर} कहते हैं।\footnote{ऊपर के नियम में $a$ !!a{constant}
है, फिर भी ऐतिहासिक कारणों से ``अभिलाक्षणिक चर'' शब्द प्रयुक्त होता है।}

\subsection{$\lexists$ के नियम}

\begin{defish}
\Axiom$ !A(a), \Gamma \fCenter \Delta $
\RightLabel{\LeftR{\lexists}}
\UnaryInf$ \lexists[x][!A(x)], \Gamma \fCenter \Delta$
\DisplayProof
\hfill
\Axiom$ \Gamma \fCenter \Delta, !A(t) $
\RightLabel{\RightR{\lexists}}
\UnaryInf$ \Gamma \fCenter \Delta, \lexists[x][!A(x)]$
\DisplayProof
\end{defish}

यहाँ भी $t$~बंद पद है और $a$~ऐसा !!a{constant} है जो \LeftR{\lexists} नियम
के निचले अनुक्रम में नहीं आता। $a$ को \LeftR{\lexists} अनुमान का
\emph{अभिलाक्षणिक चर} कहते हैं।

\RightR{\lforall} या \LeftR{\lexists} अनुमान के निचले अनुक्रम में
अभिलाक्षणिक चर के न आने की शर्त को \emph{अभिलाक्षणिक-चर शर्त} कहते हैं।

\begin{explain}
स्मरण करें: जब $!A$ ऐसा !!a{formula} हो जिसमें !!{variable}~$x$ मुक्त है, तो
इसे~$!A(x)$ लिखकर दिखाते हैं। उसी संदर्भ में $!A(t)$,
संक्षेप में~$\Subst{!A}{t}{x}$ है। अतः \RightR{\lexists} नियम को ऐसे भी लिख
सकते हैं:
\begin{prooftree}
  \Axiom$\Gamma \fCenter \Delta, \Subst{!A}{t}{x}$
  \RightLabel{\RightR{\lexists}}
  \UnaryInf$\Gamma \fCenter \Delta, \lexists[x][!A]$
\end{prooftree}
ध्यान दें कि $t$,~$!A$ में पहले से आ सकता है; जैसे $!A$ का रूप
$\Atom{\Obj P}{t,x}$ हो सकता है। इसलिए $\Gamma \Sequent \Delta,
\lexists[x][\Atom{\Obj P}{t,x}]$ को~$\Gamma \Sequent \Delta,
\Atom{\Obj P}{t,t}$ से निष्कर्षित करना
\RightR{\lexists} का सही प्रयोग है---पूर्वाधार में~$t$ की एक या अधिक, पर
आवश्यक नहीं कि सभी, आवृत्तियों को आबद्ध !!{variable}~$x$ से ``बदल'' सकते हैं।
परंतु \RightR{\lforall} और~\LeftR{\lexists} की अभिलाक्षणिक-चर शर्तें माँगती
हैं कि !!{constant}~$a$,~$!A$ में न आए। अतः
$\Gamma \Sequent \Delta, \lforall[x][\Atom{\Obj P}{a,x}]$ को
$\Gamma \Sequent \Delta, \Atom{\Obj P}{a,a}$ से~$\RightR{\lforall}$ द्वारा
सही ढंग से निष्कर्षित नहीं कर सकते।
\end{explain}

\begin{explain}
\RightR{\lexists} और \LeftR{\lforall} में पद~$t$ पर कोई प्रतिबंध नहीं है।
दूसरी ओर, \LeftR{\lexists} और \RightR{\lforall} नियमों की अभिलाक्षणिक-चर
शर्त माँगती है कि ऊपरी अनुक्रम में !!{constant}~$a$,~$!A(a)$ के बाहर कहीं न
आए। तंत्र की सुदृढ़ता---अर्थात् केवल वैध अनुक्रमों को !!{derive} करना---सुनिश्चित
करने के लिए यह आवश्यक है। इस शर्त के बिना निम्नलिखित अनुमत होता:
\begin{prooftree}
  \Axiom$!A(a) \fCenter !A(a)$
  \RightLabel{*\LeftR{\lexists}}
  \UnaryInf$\lexists[x][!A(x)] \fCenter !A(a)$
  \RightLabel{\RightR{\lforall}}
  \UnaryInf$\lexists[x][!A(x)] \fCenter \lforall[x][!A(x)]$
  \DisplayProof\bottomAlignProof
  \qquad
  \Axiom$!A(a) \fCenter !A(a)$
  \RightLabel{*\RightR{\lforall}}
  \UnaryInf$!A(a) \fCenter \lforall[x][!A(x)]$
  \RightLabel{\LeftR{\lexists}}
  \UnaryInf$\lexists[x][!A(x)] \fCenter \lforall[x][!A(x)]$
\end{prooftree}
परंतु $\lexists[x][!A(x)] \Sequent \lforall[x][!A(x)]$ वैध नहीं है।
\end{explain}

\end{document}
